% !TEX root = ../main.tex
\section{Discussion}

\hypertarget{what-drives-the-inconsistency}{%
\subsection{Drivers of inconsistency}\label{what-drives-the-inconsistency}}

The reproduction cleanly separates the ingredients of the dynamic-inconsistency phenomenon. The harvest-flow constraint is the operative between-period link. Without this constraint, the announced plan's tail remains optimal at each replanning cycle. with the constraint, however, the re-solving planner either improves upon the announced plan's tail or finds it infeasible from the realized state. The non-declining-yield policy, the most restrictive, is the most frequently inconsistent. The disequilibrium forest structure provides the underlying conditions for this inconsistency. The mature/old-growth structure creates conditions for inconsistency, while the inclusion and subsequent abandonment of negatively valued CM-CE strata provide an observable marker. The extensions further isolate the underlying mechanism (Section 4, Table~\ref{tab:extensions}). E3 shows that negatively valued CM-CE strata are a marker rather than a necessary condition for dynamic inconsistency: under a revenue-denominated floor, CM-CE harvest is structurally excluded from the projected solution, yet the announced plan is still not followed under replanning. E4 shows that the result is also robust to the specific implementation of the non-declining-yield constraint. Replacing pointwise NDY with a rolling-mean floor produces essentially the same divergence under within-plan re-anchoring, whereas anchoring the constraint to realized harvest substantially reduces the divergence, albeit at the cost of frequent constraint relaxations to maintain feasibility. Together, these results identify the interaction between a binding inter-period flow constraint, the disequilibrium initial forest structure, and its treatment under sequential replanning as the central mechanism underlying the observed dynamic inconsistency.

\hypertarget{the-discount-rate-nuance}{%
\subsection{Role of discount rates}\label{the-discount-rate-nuance}}

One finding deserves emphasis because it is easy to misstate. A plausible initial hypothesis is that the interaction between a discounted objective and an even-flow constraint gives rise to the violation of Bellman's principle, which would imply a central role of the discount rate and suggest that lowering it could restore consistency. However, the reproduction scenarios show that dynamic inconsistency occurs at every discount rate (0\%, 2\%, 4\%, and 6\%) for the flow-constrained policies. The resolution is in the thesis itself \citep[ch.~3, p.~58--59]{daugherty1991credibility}: because the discount factor is uniform across periods, it cancels out of the conditions that a plan must satisfy to be consistent. The \emph{occurrence} of inconsistency does not require a particular rate. Nonetheless, discounting shapes \emph{which} plan is chosen and the \emph{magnitude} of the divergence. The mean relative divergence and total-volume shortfall are markedly larger at a 0\% rate than at 2--6\%. The objective-gap diagnostic reveals an important distinction  that is not captured by the trajectory-divergence metric alone: at a 0\% rate the flat objective admits near-tie optima, so the trajectory-divergence metric also flags the unconstrained control, while the objective-gap diagnostic shows that its announced tail is still optimal. At positive rates, the objective discriminates. The flow-constraint divergence is genuine. So, the rate matters for the plan and for how badly the plan is diverging from the open-loop solution, but not for whether the plan is inconsistent. This result reproduces rather than contradicts the thesis, a subtle implication of the thesis: recovering the counter-intuitive rate-independence on independent software implementation provides evidence that the reproduction captures the menchanisms identified in the original analysis. To conclude, changing zhe discount-rate policy is not a remedy for non-credible plans. Extension E1 strengthens this conclusion while revealing an additional source of dynamic inconsistency. The \emph{level} of a constant rate does not govern occurrence of dynamic inconsistency, nor does a declining \emph{shape} restore credibility -- occurrence rises to 98--100\% under declining schedules. Moreover, the declining schedules introduce a preference-level time inconsistency as described by \citet{strotz1956myopia}: because each replanning planner re-applies the declining descount schedule from her own present, even the unconstrained-flow control -- which carries no structural coupling -- becomes dynamically inconsistent. The unconstrained scenarios thus distinguish between two mechanisms. Under constant 4--6\% rates, divergence arises only in the presence of a flow constraint (structural channel), whereas under declining discount schedules it also occurs in its absence. Discounting policies therefore shapes \emph{which} plan is announced and \emph{how} it diverges under replanning, but neither lowering a constant rate nor adopting a declining discount schedule restores plan credibility.

\hypertarget{implications-for-the-literature}{%
\subsection{Implications for the literature}\label{implications-for-the-literature}}

These findings have direct implications for open-loop LP forest-planning studies and planning practice rely on single-solve strategic models with inter-period harvest-flow constraints. Where such a model is applied to a disequilibrium initial forest, projected harvest trajectories, residual growing-stock trajectories, and associated shadow prices may describe a plan that is not maintained under sequential re-optimization. Trade-off analysis based on these projections may consequently differ from those obtained under a more realistic sequential implementation. The ``declining non-declining yield'' phenomenon \citep[pp.~331--332]{gunn2007models} captures this discrepancy between the harvest trajectory projected under the initial plan and that realized under sequential replanning. The same mechanism surfaces downstream in the hierarchical link between strategic and operational planning. \citet{paradis2013risk} showed that incoherent linkage between long-term wood-supply allocation and short-term harvest planning can generate systematic divergence between planned and implemented harvest. The dynamic inconsistency reproduced here provides a theoretical mechanism through which such divergence can arise and links this planning problem to the broader economics literature on even-flow constraints and the allowable-cut effect \citep{thompson1966traditional,schweitzer1972allowable,luckert1995allowable,hegan2000economic}. Likely because the original analysis remained confined to an archival thesis rather than a peer-reviewed journal publication, the mechanism has received limited direct attention in the subsequent literature. The open and reproducible implementation provided here therefore makes the original result more readily accessible for evaluation and application in forest-planning research.

\hypertarget{remedies}{%
\subsection{Remedies}\label{remedies}}

The thesis and the subsequent literature point to three families of remedies, whose implications are made more explicit by our reproduction. First, \emph{consistent formulations}: rather than solving a single open-loop problem, the optimization problem can be formulated to account explicitly for sequential decision-making, yielding subgame-perfect or feedback solutions in which each period's decision is optimal given the state and the continuation policy. The thesis already discusses the generation of consistent solutions. Their computation for the full case study represents a natural follow-on extension of the present framework. Second, \emph{credible precommitment}: inconsistency arises because future planners are free to re-solve. Institutional commitments that effectively constrain future decisions, for example, regeneration mandates attached to harvest, change the future optimization problem and may render announced plans dynamically consistent. The principal-agent literature on incentive-compatible timber-supply policy after disturbance develops this perspective further \citep{bogle2015protecting}. Third, \emph{receding-horizon practice with consistent reporting}: if planning is implemented through sequential replanning, projected outcomes should account for this decision process. The sequential-replanning simulator used here provides such a representation, and plan documents could therefore report replanned trajectories alongside, or instead of, the initial open-loop projection.

The extensions identify a fourth operational remedy and qualify the role of institutional precommitment. E2's maximum harvest-level search replaces the inter-period flow constraint with a calibrated per-period harvest cap and thereby avoids the source of dynamic coupling identified above. It identifies the highest harvest level that is both even and sustainable under sequential replanning, eliminating inconsistency in the present experiment at a level approximately 8\% below the NDY plan's initially projected harvest level. The resulting difference quantifies the reduction in harvest required to obtain an even-flow trajectory that remains feasible under replanning and provides an operational interpretation of the allowable-cut effect in this setting. Conversely, E4's realized-history formulation shows that the effectiveness of binding institutional rules remains conditional on the evolving forest state: anchoring the flow constraint to realized harvest mitigates divergence but requires relaxation in approximately one fifth to a third of replanning periods.  Thus, stronger intertemporal commitment can improve plan consistency, but cannot guarantee feasibility when the realized forest state is unable to support the prescribed harvest floor.

\hypertarget{limitations}{%
\subsection{Limitations}\label{limitations}}

Several limitations define the scope of this contribution's claims. First, the data vintage: the case study is built from the 1990 FEIS (the FORPLAN model lineage with updated 1980s prices), cross-checked against the 1987 DEIS that Daugherty used directly; the two vintages' yield tables agree, but the thesis's exact base-case schedule is not matched period by period because the thesis's per-period realized schedule is not available. Validation therefore relies on the thesis's printed anchors: the mature-type PNVs match Table~5.4 exactly (their volumes are derived from the anchors, cross-checked against the independent FEIS standing-volume curves), and the managed prescriptions match Table~5.3 in \emph{sign} and broad rotation range -- although the exact optimal-rotation ordering is only partially reproduced because the FEIS prescription mapping simplifies some treatments. Second, the formulation caveat: the thesis used a Model~II formulation; the reproduction uses Model~I because the Model~II flow-constraint compilation is not yet implemented in \texttt{ws3} -- a necessary substitution whose behavioral equivalence under sequential replanning remains to be explicitly demonstrated. Third, the thesis used ending-period (terminal) constraints with its flow-constrained runs to minimize horizon effects; we implemented these in the stack (currently opt-in and experimental) and found that with the corrected case-study data, the model does not exhibit the horizon-end liquidation these constraints are intended to prevent, so they are omitted from the reported scenarios. Fourth, the simulator implements a rolling-horizon (receding-horizon) institution. A shrinking-horizon variant (the pure Bellman tail over a fixed terminal date) shows the same qualitative phenomenon at somewhat smaller magnitude (on landbase 1 under NDY at 4\% interest rate, mean relative divergence 0.050 shrinking vs 0.074 rolling; both inconsistent). Fifth, scope: this is a single case study. The pervasiveness result (73\% of flow-constrained cells) is a property of this scenario grid, rather than an estimate of prevalence. An extension to other forests, formulations, and consistent-solution constructs is future work for which the open implementation provides a reproducible basis. Sixth, the extension grids carry their own caveats: they are additive experiments against the unchanged core control, not re-estimations of it; the E2 calibrated caps depend mildly on the recorded even-flow thresholds (a sensitivity analysis is a natural follow-on); E4's rolling-mean rows are built on pinned-\texttt{ws3} constraint internals (with a regression test degenerating them exactly to NDY); and E1 indexes discount paths by each subproblem's own present, so a calendar-absolute reading of declining schedules (which would remove the preference-drift channel while leaving the structural one) remains untested.
